\documentclass[12pt]{article}

% House style preamble
\input{.house-style/preamble.tex}

% Update PDF metadata
\hypersetup{
  pdftitle={English Passive Voice as HPC}
}

\title{English Passive Voice as HPC}
\author{Brett Reynolds \orcidlink{0000-0003-0073-7195}%
\thanks{Contact: \href{mailto:brett.reynolds@humber.ca}{brett.reynolds@humber.ca}}\\
Humber Polytechnic \& University of Toronto}
\date{\today}

\begin{document}
\maketitle

\begin{abstract}
TODO: Write abstract here.
\end{abstract}

\section{Introduction}

TODO: Write introduction here.

\section{Background}

TODO: Add background section.

\section{Analysis}

TODO: Add analysis section.

\section{Discussion}

TODO: Add discussion.

\section{Conclusion}

TODO: Write conclusion.

\section*{Acknowledgements}

% TODO: Replace with the actual models used on this project and their current versions.
% Check model names: frontier models change frequently (e.g., GPT-4 → GPT-4o → o3).
This paper was drafted with the assistance of large language models. All content has been reviewed and revised by the author, who takes full responsibility for the final text.

\newpage
\appendix

\section{From mixed cues to form/meaning split}\label{app:cue-restructuring}

The annotation scheme used in the confirmatory study went through a substantial revision between the boundary pilot and the confirmatory round. This appendix documents what the original scheme looked like, what the pilot exposed, and how the revised scheme was motivated.

\subsection{The original scheme}

The preregistered cue set contained seven coding columns plus a derived \term{peripheral subtype} label:

\begin{enumerate}
\item \term{auxiliary type} (be / get / none-other)
\item \term{participial predicate} (yes / no)
\item \term{agent realization} (by-phrase / other-overt-agentive / none)
\item \term{promotion type} (direct-object-promotion / oblique-stranding-promotion / none-unclear)
\item \term{eventive vs.\ stative} (eventive / stative / ambiguous)
\item \term{syntactic environment} (finite-clause / nonfinite-or-reduced)
\item \term{subject role profile} (patient-theme-like / locative-oblique-like / unclear-other)
\end{enumerate}

These cues were intended to capture the properties that cluster in prototypical passives. The scheme was designed for a Bayesian cue-bundle model that would test whether this cluster projects across corpora better than deterministic checklist baselines.

\subsection{What the pilot exposed}

A 16-item boundary pilot was annotated by four independent coders: the author, a base annotation, and two independent LLM agents (both Claude Opus~4.6, one with prior context and one naive). Each item was coded on 9 columns, yielding 144 cells per pair.

\begin{table}[h]
\centering
\caption{Pairwise raw agreement on 16 boundary-pilot items (144 cells per pair)}
\label{tab:pilot-agreement}
\begin{tabular}{lcc}
\hline
Pair & Disagreements & Agreement \\
\hline
Opus-2 (naive) vs.\ Base & 4 & 97.2\% \\
Opus-1 vs.\ Opus-2 (naive) & 8 & 94.4\% \\
Opus-1 vs.\ Base & 8 & 94.4\% \\
Author vs.\ Base & 15 & 89.6\% \\
Author vs.\ Opus-2 (naive) & 15 & 89.6\% \\
Author vs.\ Opus-1 & 16 & 88.9\% \\
\hline
\end{tabular}
\end{table}

The disagreements were not random. They clustered by column, as Table~\ref{tab:column-disagreements} shows.

\begin{table}[h]
\centering
\caption{Disagreements by column across all six coder pairs (96 cells per column: 16 items $\times$ 6 pairs)}
\label{tab:column-disagreements}
\begin{tabular}{lccl}
\hline
Column & Level & Disagreements & Pattern \\
\hline
\term{auxiliary type} & morphosyntactic & 3 & one item (BP009) \\
\term{participial predicate} & morphological & 0 & perfect agreement \\
\term{agent realization} & mixed & 0 & perfect agreement \\
\term{promotion type} & mixed & 14 & systematic \\
\term{eventive vs.\ stative} & semantic & 12 & systematic \\
\term{syntactic environment} & syntactic & 8 & two items \\
\term{subject role profile} & semantic & 6 & three items \\
\term{family status} & derived & 11 & cascading \\
\term{peripheral subtype} & mixed & 12 & systematic \\
\hline
\end{tabular}
\end{table}

The morphosyntactic columns (\term{auxiliary type}, \term{participial predicate}, \term{agent realization}) achieved near-perfect agreement. The semantic and mixed-level columns (\term{promotion type}, \term{eventive vs.\ stative}, \term{peripheral subtype}) produced systematic disagreements across multiple items and coder pairs. The \term{family status} disagreements were cascading: they followed from upstream disagreements in the cue columns rather than from independent judgment differences about category membership.

Two systematic patterns accounted for most of the variance. First, the author coded \term{promotion type} = none-unclear in five items where the other three coders coded direct-object-promotion. This reflected an interpretive question: do prenominal modifiers and reduced clauses without overt subjects count as having a promoted participant? The column mixed a syntactic question (is there a subject position?) with a semantic one (is there an undergoer?), and coders resolved the mix differently. Second, two items (\mention{do not be put off}, \mention{do you get charged}) produced disagreements on \term{syntactic environment} because the finite/non-finite distinction was ambiguous for imperatives and do-support constructions.

The root cause was that the original columns conflated linguistic levels. \term{Peripheral subtype} mixed morphosyntactic criteria (\mention{get}), syntactic criteria (\mention{reduced\_embedded}, \mention{prepositional}), and semantic criteria (\mention{stative\_adjectival}) in a single column. \term{Agent realization} mixed a form question (is there a \mention{by}-PP?) with a semantic one (is it agentive?). \term{Eventive vs.\ stative} was ostensibly semantic, but coders reached for syntactic diagnostics to answer it. When form and meaning came apart~-- exactly the boundary cases the study targets~-- coders had to resolve the conflict within a single cell. This is where the disagreements appeared.

\subsection{The middle-construction test}

The restructuring was motivated by a simple thought experiment. Consider the middle construction: \mention{The butter spreads easily}. Under the original scheme, three of seven cues could not distinguish this from a passive: the subject corresponds to a direct-object promotion, the subject role is patient-theme-like, and the syntactic environment is a finite clause. Only the morphosyntactic cues (no participle, no auxiliary) separated the two.

This showed that \term{promotion type}, \term{subject role profile}, and \term{syntactic environment} were tracking voice alternations generally, not the passive specifically. Middles, unaccusatives (\mention{The door opened}; \mention{Twenty people died}), tough-constructions (\mention{This book is easy to read}), and needs-washing constructions (\mention{The house needs painting}) all share the passive's functional properties~-- patient promotion to subject, agent suppression, topic management~-- while using entirely different formal means.

The conclusion was that the form cues define the property \emph{cluster} (what makes something specifically passive), while the meaning and information-structure properties explain why the cluster is \emph{maintained} (the functional pressures that keep it in use). This is the HPC architecture: form properties are the cluster; functional properties are the homeostatic mechanisms.

\subsection{The restructured scheme}

The revised scheme separates cues into level-pure layers.

The form cues are six observable morphosyntactic and constructional properties:

\begin{enumerate}
\item \term{participial form}: past participle, gerund-participial, or none
\item \term{licensing marker}: be, passive \mention{get}, causative \mention{get}/\mention{have}, subordinator or adjunct, modifier position, or absent
\item \term{constructional environment}: clausal predication, bare infinitival complement, to-infinitival complement, gerund-participial clause, adjunct participial clause, object predicative complement, or reduced modifier
\item \term{local subject present}: yes or no
\item \term{by-PP present}: yes or no
\item \term{stranded preposition}: yes or no
\end{enumerate}

The meaning cues are three judgment-based semantic properties of the predicate, not the sentence:

\begin{enumerate}
\setcounter{enumi}{6}
\item \term{event implied}: does this predicate imply an event occurred?
\item \term{agent implied}: is an agent or causer understood?
\item \term{predicand as undergoer}: is the predicand the affected participant?
\end{enumerate}

The key structural change is that \term{constructional environment} replaces \term{syntactic environment} without invoking a finite/non-finite binary. It names construction types directly (clausal predication, bare infinitival complement, to-infinitival complement, etc.) rather than forcing them through a binary clause-type split. This resolved the pilot's most persistent disagreement: whether imperatives (\mention{do not be put off}) and do-support constructions (\mention{do you get charged}) should count as finite. Under the revised scheme, both are bare infinitival complements, and the question does not arise.

\subsection{The comparator correction}

The pilot exposed a second design problem, this time in the deterministic
baselines rather than in the annotation columns. An initial v2 draft defined
the strict checklist over the annotated field \term{licensing marker}. That was
too close to the analyst's own structural judgment: the checklist was no longer
a genuinely surface-oriented competitor, but a partial recoding of the gold
analysis. On the corrected v2 fit-gate pilot, that version of the checklist was
perfect on the observed \term{core}/\term{foil} rows, and the fake-data gate
failed.

The preregistered correction was to redefine the checklists over overt surface
\mention{be}/\mention{get} material in the retained extraction metadata, plus a
small number of annotated convenience cues. This matters methodologically. The
goal of the comparison is not to let the Bayesian model beat a straw man, but
it is also not to let the checklist inherit the analyst's own functional
judgment. The corrected baseline is therefore strong, but still capable of
overgenerating on passive-looking foils such as copular participial clauses.

The first corrected fake-data run, at the minimum preregistered resolution of
200 simulations, landed at 0.745 on the strict gate and 0.835 on the stronger
gate. Because that strict result missed the threshold by one simulated pass, it
was treated as a borderline Monte Carlo estimate rather than as a substantive
design failure. A preregistered higher-resolution rerun at 1000 simulations
yielded 0.788 and 0.844, respectively, which is the estimate used for the
final adequacy decision.

\subsection{Comparison constructions}

The revised item set includes tokens from non-passive constructions that share meaning and information-structure properties with passives: middles, unaccusatives, tough-constructions, needs-washing constructions, and \mention{get} + adjective. These receive \term{family status} = foil and are coded on all cues. They test whether the form cluster holds together even when the meaning cues overlap with non-passives.

The needs-washing construction is where \term{participial form} = gerund-participial does its work. \mention{The house needs painting} is participial, but it is the wrong participle for passive. A binary participial-predicate column would miss this distinction.

\subsection{What the revision shows}

The original scheme treated the passive as a single cluster to be coded at a single level. The revised scheme treats it as a cluster defined at the form level and maintained by functional pressures at the meaning and information-structure levels. The boundary cases~-- the theoretical target of the study~-- are precisely the cases where form and meaning come apart. By separating the levels, the scheme lets the model see the dissociation rather than forcing the coder to resolve it.

\newpage
\printbibliography

\end{document}
